\section{Magnetic domains and joint-prime gluing}\label{sec:gluing}
\subsection{What higher magnetic powers would require}
The domain statement for one insertion is sufficient for
Theorem~\ref{thm:schur-prime}. It does not place the prepared state in
the domain of every negative magnetic power.
\begin{proposition}[Formal higher-power coefficient criterion]
\label{prop:magnetic-powers}
For $k\geq1$, iteration of the coefficient action gives the formal series
\begin{equation}\label{eq:magnetic-k}
 (u_-^k\eta_a)_{-k}(\zeta)=\sqrt{1-a^2}
 \frac{\prod_{j=0}^{k-1}(1+q^{k-1-2j}\zeta)}{1+aq^{-k}\zeta}.
\end{equation}
Its Taylor coefficient sequence belongs to $\ell^2$ if and only if
\begin{equation}\label{eq:magnetic-criterion}
 a<q^k\quad\hbox{or}\quad a\in\{q,q^3,\ldots,q^{2k-1}\}.
\end{equation}
For fixed $0<a<1$, square summability of all these formal outputs is
equivalent to $a=q^{2j+1}$ for some integer $j\geq0$.
\end{proposition}
\begin{proof}
Induction in \eqref{eq:schur-actions} gives \eqref{eq:magnetic-k}.
For every coefficient of degree $n\geq k$, expanding the denominator
shows that the coefficient equals
\begin{equation}\label{eq:magnetic-tail}
 \sqrt{1-a^2}(-aq^{-k})^n
 \prod_{j=0}^{k-1}(1-q^{2j+1}/a).
\end{equation}
If the product is nonzero, the geometric tail is square summable
exactly when $aq^{-k}<1$. If the product vanishes, the numerator
cancels the denominator and the output is a polynomial. This proves
\eqref{eq:magnetic-criterion}. For all $k$, the inequality eventually
fails at any fixed $a>0$, so an exceptional value is necessary.
If $a=q^{2j+1}$, it satisfies the strict inequality for $k\leq j$
and the cancellation condition for $k\geq j+1$, proving sufficiency
for the formal coefficient statement.
\end{proof}

The exceptional square-summable outputs do not by themselves prove
membership in powers of a chosen closed operator extension. A divergent
coefficient tail is, however, a necessary domain obstruction whenever
the prescribed action on Fourier coefficients is retained. This
distinction prevents meromorphic continuation across a pole from being
mistaken for an operator-domain argument.

For the prime states $a_p=qp^{-1/2}$, requiring all negative powers
would force $p^{-1/2}=q^{2j_p}$. This cannot hold simultaneously for
$p=2$ and $p=3$, since it would imply a nontrivial equality between
an integer power of $2$ and an integer power of $3$. It does not affect
the single insertion already proved. More generally, for any fixed
finite $K>1$, choosing $q^{K-1}>2^{-1/2}$ puts every prime state in
the strictly decaying regime for all powers through $K$. In that
regime the Taylor-approximation argument of
Lemma~\ref{lem:one-magnetic} extends through the finite sequence of
shifts. A proposed junction must specify the finite words or actual
operator domains it uses; a field theory need not put every state in
the domain of every unbounded power.

\subsection{A two-prime obstruction for fixed output vectors}
For each prime write $r_p=p^{-1/2}$, $d_p=\log p$, and
\begin{equation}\label{eq:prime-filter}
 A_p(z)=\frac{1-z}{1-r_pz}
 =1-(1-r_p)\sum_{n\geq1}r_p^{n-1}z^n,
 \qquad c_p=\frac{d_pr_p(1+r_p)}{1-r_p}.
\end{equation}
The norm of $\sqrt{c_p}A_p$ is the multiplier
\eqref{eq:prime-weight}. Consider a proposed coherent sum with fixed
unit output vectors $v_2,v_3$:
\begin{equation}\label{eq:constant-vector-source}
 F\longmapsto\widehat F(\tau)
 \left[\sqrt{c_2}A_2(e^{id_2\tau})v_2
       +\sqrt{c_3}A_3(e^{id_3\tau})v_3\right],
\end{equation}
with output measure $\dd\tau/(2\pi)$.

\begin{proposition}[Mixed atom from constant-vector gluing]\label{prop:mixed}
In \eqref{eq:constant-vector-source}, the coefficient of
$e^{i(\log3-\log2)\tau}$ in the full norm multiplier is
\begin{equation}\label{eq:mixed-atom}
 \sqrt{c_2c_3}(1-r_2)(1-r_3)\ip{v_2}{v_3}.
\end{equation}
Consequently this source, together with independently added gamma,
contact and pole terms, cannot equal the two-prime Weil form on
$I_{5/4}$ unless $v_2\perp v_3$. In the orthogonal case its prime
part is exactly the sum of the separate positive references.
\end{proposition}
\begin{proof}
The first cross term in the squared norm is
$\sqrt{c_2c_3}\overline{A_2(e^{id_2\tau})}
A_3(e^{id_3\tau})\ip{v_2}{v_3}$.
Its first nonconstant coefficient from each factor gives
\eqref{eq:mixed-atom}. A second pair of nonnegative indices with
$n\log3-m\log2=\log3-\log2$ would give
$3^{n-1}=2^{m-1}$ and hence $n=m=1$. The reverse cross term cannot
give this positive frequency with nonnegative indices, and the
individual squared norms contain only integer multiples of $d_2$
or $d_3$.

The Fourier coefficient series of each factor is absolutely summable;
the cross kernel is therefore a finite-variation atomic measure. Other
mixed lengths may accumulate, but its atom at $\log(3/2)$ has the
specified nonzero mass if the overlap is nonzero. The Weil target has
no prime-power atom there. Its gamma and pole kernels are smooth near
this nonzero separation, and its contact is supported at zero, so none
can cancel the atom. Since $0<\log(3/2)<5/4$, restriction to $I_{5/4}$
still detects it. Equality of all polarized test-function pairings
would imply equality of their distribution kernels, a contradiction.
A purely imaginary overlap gives an imaginary antisymmetric kernel;
complex tests still detect it. Thus the overlap must vanish.
\end{proof}

The proposition does not apply to arbitrary coherent gluing. A joint
operator-valued junction may have additional terms which cancel mixed
returns as an identity. Producing and evaluating such a junction is a
new source construction. The failed fixed-vector ansatz gives no reason
to estimate its unwanted atoms as a small residual.

\subsection{The full comparison, with no omitted contact or pole}
Let $\mathcal P$ be any finite set containing all primes $p<e^L$.
Equations~\eqref{eq:target} and \eqref{eq:lifted-norm} give the exact
identity
\begin{align}
 Q_L[f]={}&K[E_Lf]+\sum_{p\in\mathcal P}B_{r_p,d_p}[E_Lf]
 \notag\\
 &+\left(w_0-\sum_{p\in\mathcal P}\kappa_{r_p,d_p}\right)\norm f^2
 +P_L[f].
 \label{eq:full-comparison}
\end{align}
Each term in the first line is independently positive. The second
line has not been realized as part of their joint norm. Extra primes
in $\mathcal P$ contribute only contacts on this interval, which
cancel algebraically against the displayed subtraction. This does
not make the subtraction an independently positive operation.

Thus the dressed result arrives at the same contact and pole problem
as the conservative prime reference. It supplies a concrete Schur
preparation and removes the earlier undressed-state mismatch, but
establishes no larger positivity interval. Unitary repackaging of the
orthogonal gamma and prime outputs leaves \eqref{eq:full-comparison}
unchanged. The next calculation must determine contact, poles and
prime interference within one complete positive preparation.

\subsection{Interference is necessary, and the two-channel architecture fails}
\label{subsec:interference}
Equation~\eqref{eq:full-comparison} presents $Q_L$ as a sum of two
independently positive forms and a residual. Write
\begin{equation}\label{eq:channels}
 A_L[f]=K[E_Lf],\qquad
 B_L[f]=\sum_{p\in\mathcal P}B_{r_p,d_p}[E_Lf],\qquad
 R_L[f]=\Big(w_0-\sum_{p\in\mathcal P}\kappa_{r_p,d_p}\Big)\norm f^2+P_L[f],
\end{equation}
so that $Q_L=A_L+B_L+R_L$. The residual is not a positive form, and the
first observation is that it cannot be repaired by adding more positive
channels.

\begin{proposition}[No orthogonal augmentation]\label{prop:no-augment}
Let $c<0$. On the subspace
$\{f\in C_c^\infty(I_L):\int_{I_L}f\cosh(x/2)=\int_{I_L}f\sinh(x/2)=0\}$,
which has codimension two, the form $c\norm f^2+P_L[f]$ equals
$c\norm f^2<0$. Hence $R_L$ is not the Gram form of any further channel
adjoined orthogonally to a preparation whose norm is $A_L+B_L$.
\end{proposition}
\begin{proof}
$P_L$ vanishes on that subspace by \eqref{eq:poles}, and the subspace is the
common kernel of two linear functionals on an infinite-dimensional space.
\end{proof}

The remaining possibility is interference: the archimedean and arithmetic
parts of a preparation must fail to be orthogonal. The next elementary bound
says how much interference is available.

\begin{lemma}[Interference bound]\label{lem:interference}
Let $G,S$ be complex-linear maps from a complex vector space $V$ into a
Hilbert space $\HH$. Put $A[f]=\norm{Gf}^2$, $B[f]=\norm{Sf}^2$ and
$C[f]=\norm{(G+S)f}^2-A[f]-B[f]$. Then
\begin{equation}\label{eq:interference-bound}
 tA[f]+t^{-1}B[f]+C[f]\geq0\qquad (t>0,\ f\in V),
\end{equation}
equivalently $-C[f]\leq2\sqrt{A[f]B[f]}$, that is
$\norm{(G+S)f}^2\geq\big(\sqrt{A[f]}-\sqrt{B[f]}\big)^2$.
\end{lemma}
\begin{proof}
$C[f]=2\re\ip{Gf}{Sf}$, so
$tA[f]+t^{-1}B[f]+C[f]=\norm{\sqrt t\,Gf+t^{-1/2}Sf}^2\geq0$.
Minimizing the left side over $t>0$ at $t=\sqrt{B/A}$ gives the other two
forms, and replacing $S$ by $-S$ leaves the statement unchanged.
\end{proof}

At $t=1$ the bound \eqref{eq:interference-bound} is only the positivity of
$Q_L$. Other values of $t$ are strictly stronger, and they exclude the
architecture of \eqref{eq:full-comparison}.

\begin{theorem}[Exclusion of the archimedean/prime two-channel source]
\label{thm:two-channel}
Let $L=1$, or $L=5/4$, or let $L$ be sufficiently large. There are no
complex-linear maps $G,S$ from $C_c^\infty(I_L)$ into a common positive
Hilbert space with
\begin{equation}\label{eq:two-channel-hypothesis}
 \norm{Gf}^2=A_L[f],\qquad \norm{Sf}^2=B_L[f],\qquad
 \norm{(G\pm S)f}^2=Q_L[f]
\end{equation}
for every $f\in C_c^\infty(I_L)$.
\end{theorem}
\begin{proof}
For such a pair the interference form is $C=Q_L-A_L-B_L=R_L$, so
Lemma~\ref{lem:interference} requires
\begin{equation}\label{eq:two-channel-test}
 N_L(t)[f]:=tA_L[f]+t^{-1}B_L[f]+R_L[f]\geq0
\end{equation}
for every $t>0$. Each of $B_L$, $P_L$ and $\norm\cdot^2$ is bounded and
$A_L$ is the closed gamma energy, so $N_L(t)$ is continuous for the energy
norm of $\mathcal D_{\gamma,L}$. Moreover $C_c^\infty(I_L)$ is dense in
$\mathcal D_{\gamma,L}$: if $g$ is bounded by one and supported in a set of
measure $2\varepsilon$ then
$2\norm g^2-2\re\ip g{U_rg}\leq2\min(r,\varepsilon)$ pointwise in $r$, so
$K[g]=O(\varepsilon\log(1/\varepsilon))$, and cutting $\mathbf 1_{I_L}$ off
smoothly within $\varepsilon$ of each endpoint changes the energy by that
much. It therefore suffices to exhibit one $f\in\mathcal D_{\gamma,L}$ and
one $t>0$ with $N_L(t)[f]<0$.

Take $f=\mathbf 1_{I_L}$, so that $\norm f^2=L$ and
$\ip F{U_rF}=(L-r)^+$. Then
\begin{equation}\label{eq:indicator-energy}
 A_L[\mathbf 1_{I_L}]
 =\int_0^\infty\big[2L-2(L-r)^+\big]\gammak(r)\dd r
 =2\sum_{n\geq0}\frac{1-e^{-a_nL}}{a_n^2}
 \;\leq\;\tfrac12\psi'(\tfrac14),
\end{equation}
using $\int_0^Lre^{-ar}\dd r+Le^{-aL}/a=(1-e^{-aL})/a^2$ and
$\sum_{n\geq0}a_n^{-2}=\tfrac14\psi'(\tfrac14)$. Likewise
\begin{equation}\label{eq:indicator-pieces}
 B_L[\mathbf 1_{I_L}]
 =\sum_{p}\Big[\kappa_{r_p,d_p}L-2d_p\!\!\sum_{nd_p<L}\!\!r_p^{\,n}(L-nd_p)\Big],
 \qquad
 P_L[\mathbf 1_{I_L}]=32\sinh^2(L/4),
\end{equation}
the second because $\int_{I_L}\sinh(x/2)\dd x=0$ and
$\int_{I_L}\cosh(x/2)\dd x=4\sinh(L/4)$.

For $L=5/4$ the active translations are $\log2$ and $\log3$, and
\eqref{eq:indicator-energy}--\eqref{eq:indicator-pieces} give
\[
 A_L=4.302142,\quad B_L=7.197439,\quad P_L=3.228059,\quad
 \kappa_{r_2,d_2}+\kappa_{r_3,d_3}=6.348275,
\]
so $R_L=-11.422514$ and $N_L(1)=Q_L[\mathbf 1_{I_L}]=0.077067>0$, while at
$t=13/10$
\[
 N_L(13/10)=5.592785+5.536491-11.422514=-0.293238<0 .
\]
For $L=1$ the same computation with the eight-cell step function
$f=(1,1,0,-1,-1,0,1,1)$ on cells of width $1/8$ and $t=3/4$ gives
$N_L(t)[f]=-0.106621<0$, again against $Q_L[f]=0.187826>0$. Both evaluations
are finite closed forms in $\psi'(1/4)$, $\sinh$, $\log2$ and $\log3$; they
are carried out with an explicit tail bound for the gamma series in
Appendix~\ref{app:checks}.

For large $L$ take $f=\mathbf 1_{I_L}$ again. By \eqref{eq:indicator-energy}
$A_L\leq\tfrac12\psi'(\tfrac14)$ uniformly in $L$, and by
\eqref{eq:indicator-pieces} $B_L\leq L\sum_{p<e^L}\kappa_{r_p,d_p}$, since
the subtracted correlations are nonnegative. The prime number theorem and
partial summation give
$\sum_{p<x}\kappa_{r_p,d_p}=4\sqrt x\,(1+o(1))$, so
\[
 2\sqrt{A_LB_L}=O\big(L^{1/2}e^{L/4}\big),
 \qquad
 -R_L=\Big(\sum_{p<e^L}\kappa_{r_p,d_p}-w_0\Big)L-32\sinh^2(L/4)
 =4(L-2)e^{L/2}(1+o(1)).
\]
Hence $-R_L>2\sqrt{A_LB_L}$ for all large $L$, and
\eqref{eq:interference-bound} fails by a factor tending to infinity.
\end{proof}

\begin{remark}\label{rem:scan}
The failure is not confined to the two certified intervals. Evaluating
\eqref{eq:two-channel-test} at $f=\mathbf 1_{I_L}$ and $t=\sqrt{B_L/A_L}$
on a grid of step $1/10$ shows $N_L(t)<0$ at every $L\geq11/10$ up to
$L=13/2$; at $L=1$ the indicator alone gives $+0.054$ and the eight-cell
step function above is needed. These evaluations are recorded in
Appendix~\ref{app:checks}.
\end{remark}

Theorem~\ref{thm:two-channel} excludes an architecture, not the identities it
is assembled from. Theorem~\ref{thm:schur-prime} remains correct; what fails
is the plan of adjoining to it a separately positive archimedean channel and
recovering \eqref{eq:target}. The exclusion is specific to the pair
\eqref{eq:channels}: a preparation whose channel norms are different forms
$A'$, $B'$ with $A'+B'+C'=Q_L$ is not covered, and
\eqref{eq:interference-bound} is then the first test such a preparation must
pass. A source in which archimedean and prime data are not separated into two
channels at all is untouched by the argument.
Proposition~\ref{prop:mixed} and Theorem~\ref{thm:two-channel} are
complementary. The first forbids a coherent sum with fixed output vectors,
because of the mixed translations it creates. The second forbids the
orthogonal-channel bookkeeping of \eqref{eq:full-comparison} however the two
channels are subsequently correlated.

\subsection{The mirror reference and the subtraction form}
\label{subsec:subtraction}
Lemma~\ref{lem:interference} bounds the norm of a \emph{sum} from below. It
says nothing about a difference $\norm{Gf}^2-\norm{\Pi Gf}^2$, and
Proposition~\ref{prop:no-augment} says that a difference is what is needed.
This subsection shows that the target has such a presentation, with the
arithmetic confined to the subtracted term.

The positive reference of Theorem~\ref{thm:schur-prime} has the least contact
compatible with prime atoms of one sign (Corollary~\ref{cor:contact}). There
is a second elementary weight with the same atoms and the opposite sign.

\begin{proposition}[Mirror prime reference]\label{prop:mirror}
For $0<r<1$ and $d>0$ set
\begin{equation}\label{eq:mirror-weight}
 \widetilde w_{r,d}(z)=\frac{dr(1-r)}{1+r}\frac{|1+z|^2}{|1-rz|^2},
 \qquad|z|=1 .
\end{equation}
Then $\widetilde w_{r,d}\geq0$ and
\begin{equation}\label{eq:mirror-coefficients}
 \widetilde w_{r,d}(z)
 =\widetilde\kappa_{r,d}+d\sum_{n\geq1}r^n(z^n+z^{-n}),
 \qquad
 \widetilde\kappa_{r,d}=\frac{2dr}{1+r} .
\end{equation}
The constant $\widetilde\kappa_{r,d}$ is the least one for which a continuous
scalar multiplier with Fourier coefficients $+dr^{|n|}$ is nonnegative, and
\begin{equation}\label{eq:complementary}
 w_{r,d}+\widetilde w_{r,d}=\frac{4dr}{1-r^2},
 \qquad
 \widetilde\kappa_{r,d}<\kappa_{r,d}=\frac{2dr}{1-r}.
\end{equation}
\end{proposition}
\begin{proof}
Nonnegativity is clear from \eqref{eq:mirror-weight}. Write $c=z+z^{-1}$ and
use the Poisson identity \eqref{eq:poisson}. The constant Laurent coefficient
of $(2+c)/|1-rz|^2$ is $(2+2r)/(1-r^2)=2/(1-r)$, and for $m\geq1$ its $z^m$
coefficient is
\[
 \frac{2r^m+r^{m-1}+r^{m+1}}{1-r^2}
 =\frac{r^{m-1}(1+r)^2}{1-r^2}
 =\frac{r^{m-1}(1+r)}{1-r}.
\]
Multiplying by $dr(1-r)/(1+r)$ gives \eqref{eq:mirror-coefficients}.
Adding \eqref{eq:prime-weight} cancels every nonconstant coefficient and
leaves $\kappa_{r,d}+\widetilde\kappa_{r,d}=4dr/(1-r^2)$. For minimality,
a multiplier $C+d\sum_{n\geq1}r^n(z^n+z^{-n})$ takes at $z=-1$ the value
$C-2dr/(1+r)$, which must be nonnegative; \eqref{eq:mirror-weight} vanishes
there.
\end{proof}

Applying the lift \eqref{eq:bloch} to \eqref{eq:mirror-coefficients} exactly
as in Section~\ref{sec:dressed} gives a positive form on $C_c^\infty(I_L)$,
\begin{equation}\label{eq:mirror-form}
 \widetilde B_{r,d}[F]
 =\widetilde\kappa_{r,d}\norm F^2
  +2d\sum_{n\geq1}r^n\re\ip F{U_{nd}F}\geq0 .
\end{equation}

\begin{theorem}[Prime-free subtraction form of the target]\label{thm:subtraction}
Let $\mathcal P$ be any finite set of primes containing every $p<e^L$. Then
\begin{align}
 Q_L[f]&=A_L[f]-\sum_{p\in\mathcal P}\widetilde B_{r_p,d_p}[E_Lf],
 \label{eq:subtraction}\\
 A_L[f]&=K[E_Lf]
  +\Big(w_0+\sum_{p\in\mathcal P}\widetilde\kappa_{r_p,d_p}\Big)\norm f^2
  +P_L[f].
 \label{eq:prime-free}
\end{align}
The form $A_L$ involves no prime translation: away from the diagonal its
kernel is $-\gammak(|x-y|)+2\cosh((x-y)/2)$, which is smooth, and its only
arithmetic input is the constant in \eqref{eq:prime-free}. Primes outside
$e^L$ contribute only contacts, which cancel against that constant.
\end{theorem}
\begin{proof}
By \eqref{eq:mirror-form} the prime term of \eqref{eq:target} is
$-\sum_p\big(\widetilde B_{r_p,d_p}-\widetilde\kappa_{r_p,d_p}\norm f^2\big)$,
since for $F=E_Lf$ the correlations with $nd_p\geq L$ vanish and the
coefficient of $\re\ip F{U_{nd_p}F}$ is $-2d_pr_p^{\,n}$ in both. Substituting
in \eqref{eq:target} gives \eqref{eq:subtraction}. The kernel statement is
\eqref{eq:gamma-energy} and \eqref{eq:poles}.
\end{proof}

\begin{corollary}[The criterion as a domination, and as a compression]
\label{cor:domination}
For each $L>0$, $Q_L\geq0$ if and only if
$\sum_{p}\widetilde B_{r_p,d_p}\preceq A_L$ on $C_c^\infty(I_L)$. If in
addition $A_L[f]=\norm{Gf}^2$ for a complex-linear $G$ into a positive
Hilbert space and
$\sum_p\widetilde B_{r_p,d_p}[E_Lf]=\ip{Gf}{\Pi\,Gf}$ for a positive
contraction $\Pi$ on the closure of the range of $G$, then
\begin{equation}\label{eq:compression}
 Q_L[f]=\norm{(1-P)VGf}^2
\end{equation}
for an isometry $V$ and an orthogonal projection $P$ on a larger space. In
particular $Q_L$ is realized as a compression, and
Lemma~\ref{lem:interference} places no obstruction on this presentation.
\end{corollary}
\begin{proof}
The equivalence is \eqref{eq:subtraction}. Given $\Pi$, we have
$Q_L[f]=\ip{Gf}{(1-\Pi)Gf}$, and $1-\Pi$ is a positive contraction; by the
dilation theorem $1-\Pi=V^*(1-P)V$ for an isometry $V$ and a projection $P$,
which gives \eqref{eq:compression}. Lemma~\ref{lem:interference} concerns
$\norm{(G+S)f}^2$ and does not apply to \eqref{eq:compression}.
\end{proof}

Two features make \eqref{eq:subtraction} a better construction target than
\eqref{eq:full-comparison}, and one numerical observation says what is still
hard.

\emph{The mirror contact is much smaller.} Since
$\widetilde\kappa_{r,d}/\kappa_{r,d}=(1-r)/(1+r)$, the saving is concentrated
on the small primes, which is exactly where the budget of
\eqref{eq:full-comparison} was lost:
\begin{center}
\small
\begin{tabular}{lrrr}
\toprule
$L$ & $\sum_p\kappa_{r_p,d_p}$ & $\sum_p\widetilde\kappa_{r_p,d_p}$
 & $w_0+\sum_p\widetilde\kappa_{r_p,d_p}$\\
\midrule
$1$ & $3.3468$ & $0.5742$ & $-4.7980$\\
$5/4$ & $6.3483$ & $1.3785$ & $-3.9937$\\
$2$ & $11.3172$ & $3.4406$ & $-1.9315$\\
$3$ & $18.9237$ & $7.8704$ & $+2.4983$\\
$4$ & $30.3436$ & $16.0103$ & $+10.6381$\\
\bottomrule
\end{tabular}
\end{center}
The constant in \eqref{eq:prime-free} changes sign near $L=2.6$; beyond that
point the prime-free form has a positive contact.

\emph{The prime-free form is not marginal.} In a Fourier--Galerkin space of
$61$ modes on $I_L$, the least eigenvalue of $A_L$ relative to
$\norm\cdot^2$ is $0.258$ at $L=1$, $0.799$ at $L=5/4$, $0.556$ at $L=3/2$,
$2.13$ at $L=2$ and $22.4$ at $L=5$, while the same computation gives
$\lambda_{\min}(Q_L)\leq10^{-6}$ for $L\geq1$. The difficulty has moved out
of the source and into the subtraction.

\emph{The domination is saturated.} Writing
$\Pi_L=A_L^{-1/2}\big(\sum_p\widetilde B_{r_p,d_p}\big)A_L^{-1/2}$, the
criterion of Corollary~\ref{cor:domination} is $\norm{\Pi_L}\leq1$, and the
same computation returns $\norm{\Pi_L}=1$ to within $10^{-7}$ with the number
of eigenvalues above $0.9$ growing from two at $L=4/5$ to thirteen at $L=2$.
That is expected: $\norm{\Pi_L}=1-\min\operatorname{spec}(Q_L;A_L)$ and $Q_L$
is nearly degenerate. For contrast, the corresponding quotient for
\eqref{eq:full-comparison}, which asks $\sum_pB_{r_p,d_p}\preceq K$, equals
$1.75$ at $L=5/4$, $3.57$ at $L=2$ and $26.1$ at $L=5$. A subspace value
above one certifies that a domination fails, so those numbers are the
quantitative form of Theorem~\ref{thm:two-channel}; a subspace value below
one certifies nothing, so the first row is evidence and not a proof.

\begin{corollary}[Schur realization of the mirror reference]
\label{cor:mirror-schur}
Fix $0<q<1$ and $0<r<1$ and put $a=qr$. The dressing
\begin{equation}\label{eq:mirror-dressing}
 \widetilde D_{q,a}(w)=\frac{\sqrt{1-a^2}}{C_q}
 \frac{(1-qw)\,(-q^3w;q^2)_\infty}{1+aw},
\end{equation}
which is \eqref{eq:dressing} with the sign of its first Pochhammer factor
reversed, is analytic except for a simple pole at $w=-a^{-1}$, of modulus
$(qr)^{-1}>q^{-1}$. Its electric state is
$\widetilde\eta_{qr}(w)=\sqrt{1-a^2}(1-qw)/\big((1+qw)(1+aw)\big)$, one negative
magnetic insertion gives
\begin{equation}\label{eq:mirror-state}
 (u_-\widetilde\eta_{qr})_{-1}(\zeta)
 =\sqrt{1-q^2r^2}\;\frac{1-\zeta}{1+r\zeta},
\end{equation}
and with
$\widetilde\beta_{q;r,d}^{\,2}=dr(1-r)\big/\big((1+r)(1-q^2r^2)\big)$, which is
\eqref{eq:beta} with $(1+r)/(1-r)$ inverted, the preparation
$\widetilde\Phi_{q;r,d}[h]
 =\widetilde\beta_{q;r,d}\,h(-q^{-1}v)\,u_-\widetilde D_{q,a}(v)\Omega_q$
has pairing $\int_{S^1}\overline{h(z)}k(z)\widetilde w_{r,d}(z)\dd\theta/2\pi$
with $\widetilde w_{r,d}$ as in \eqref{eq:mirror-weight}. As in
Theorem~\ref{thm:schur-prime}, one fixed $q$ serves every prime.
\end{corollary}
\begin{proof}
Apply Proposition~\ref{prop:reachable} with
$\psi(\zeta)=(1-\zeta)/(1+r\zeta)$, which is analytic on $|\zeta|<r^{-1}$ and
$r^{-1}>1$. Then \eqref{eq:reach-dressing} gives
$D(w)\propto(1-qw)(-q^3w;q^2)_\infty/(1+qrw)$, which is
\eqref{eq:mirror-dressing}, and the output state is $\psi$ up to the stated
normalization. With $z=-\zeta$ one has $|1-\zeta|^2=|1+z|^2$ and
$|1+r\zeta|^2=|1-rz|^2$, and
$\widetilde\beta^{\,2}_{q;r,d}(1-q^2r^2)=dr(1-r)/(1+r)$ is the constant in
\eqref{eq:mirror-weight}.
\end{proof}

Both references are therefore states of the same representation at the same
fixed $q$, and the mirror does not require its own protected sector. The two
output states $\psi_\pm=(1\pm\zeta)/(1+r\zeta)$ are orthogonal, with
$\norm{\psi_\pm}^2=2/(1\pm r)$, and in the basis $(1+r\zeta)^{-1}$,
$\zeta(1+r\zeta)^{-1}$ their Gram matrix is
$(1-r^2)^{-1}\left(\begin{smallmatrix}1&-r\\-r&1\end{smallmatrix}\right)$,
positive definite. By Proposition~\ref{prop:reachable} the freedom exhibited
here is freedom of weights alone; a positive sector supplies no null vectors,
and Section~\ref{subsec:leaves} says why none are needed.

\subsection{The target as a compression of the gamma energy}
\label{subsec:gammacompression}
Corollary~\ref{cor:domination} presents $Q_L$ as a compression of a source for
$A_L$. That presentation requires two things not in hand: a source for $A_L$,
and the inequality $A_L\geq0$, which is \eqref{eq:prime-free-target} and is
settled only outside a compact window (Section~\ref{subsec:leaves}). Splitting
the pole form by parity removes both requirements.

\begin{proposition}[Compression of the gamma energy]
\label{prop:gamma-compression}
Let $\mathcal P$ be any finite set of primes containing every $p<e^L$, put
$c_L=w_0+\sum_{p\in\mathcal P}\widetilde\kappa_{r_p,d_p}$,
$c_+=\max(c_L,0)$ and $c_-=\max(-c_L,0)$, and define on $C_c^\infty(I_L)$
\begin{align}
 K_+[f]&=K[E_Lf]+c_+\norm f^2+2\big|\ip{\cosh(x/2)}f\big|^2,
 \label{eq:source-side}\\
 T_L[f]&=c_-\norm f^2+2\big|\ip{\sinh(x/2)}f\big|^2
        +\sum_{p\in\mathcal P}\widetilde B_{r_p,d_p}[E_Lf].
 \label{eq:subtracted-side}
\end{align}
Then $Q_L=K_+-T_L$. The form $K_+$ is a norm, $K_+[f]=\norm{Gf}^2$ for the
source
\begin{equation}\label{eq:explicit-source}
 Gf=\Big(b(D)^{1/2}E_Lf,\;c_+^{1/2}f,\;\sqrt2\,\ip{\cosh(x/2)}f\Big)
 \in L^2(\R)\oplus L^2(I_L)\oplus\C,
\end{equation}
where $b(D)^{1/2}$ is the Fourier multiplier with symbol
$b(\tau^2)^{1/2}$; and $T_L\geq0$, each of its three terms separately.
Consequently $Q_L\geq0$ on $C_c^\infty(I_L)$ if and only if
$T_L\preceq K_+$, and in that case
\begin{equation}\label{eq:gamma-compression}
 Q_L[f]=\ip{Gf}{(1-\Pi)Gf}
\end{equation}
for the positive contraction $\Pi$ on the closure of the range of $G$
determined by $\ip{Gf}{\Pi Gg}=T_L(f,g)$.
\end{proposition}
\begin{proof}
Substituting \eqref{eq:parity-split} for $P_L$ in \eqref{eq:prime-free}, and
moving the constant to the side fixed by its sign, turns
\eqref{eq:subtraction} into $Q_L=K_+-T_L$; nothing is cancelled or estimated.
$K_+$ is a norm because $b\geq0$ in \eqref{eq:gamma-multiplier}, so
$K[F]=\norm{b(D)^{1/2}F}^2$ by \eqref{eq:gamma-fourier}, while the other two
terms are a nonnegative multiple of $\norm f^2$ and a rank-one square.
$T_L\geq0$ because $c_-\geq0$, the second term is a rank-one square, and
$\widetilde B_{r_p,d_p}\geq0$ by \eqref{eq:mirror-form}. For the last
statement, $t(Gf,Gg)=T_L(f,g)$ is well defined on the range of $G$ precisely
because $T_L\preceq K_+$: if $Gf=Gf'$ then $K_+[f-f']=0$, hence
$T_L[f-f']=0$. It satisfies $0\leq t(u,u)\leq\norm u^2$ there, so the Riesz
representation supplies a positive contraction $\Pi$ with
$t(u,v)=\ip u{\Pi v}$, and \eqref{eq:gamma-compression} follows.
\end{proof}

Three differences from Corollary~\ref{cor:domination} are worth stating.

\emph{The dominating side has an explicit source.} \eqref{eq:explicit-source}
is a Fourier multiplier together with two elementary coordinates, and its
positivity is that of the digamma partial-fraction expansion
\eqref{eq:gamma-multiplier}, which is prior to any arithmetic input. That is
what distinguishes it from the observation in Section~\ref{sec:positive} that
$\Phi=Q_L^{1/2}f$ satisfies Definition~\ref{def:realization} on one interval as
soon as $Q_L\geq0$: there the positivity assumed is the statement to be proved,
here it is not. Appendix~\ref{app:gamma} realizes the same gamma form as a
minimized positive field energy, and a candidate model is being asked to
replace \eqref{eq:explicit-source} by a physically defined preparation rather
than to supply positivity.

\emph{No unproved inequality enters the presentation.}
Corollary~\ref{cor:domination} needs $A_L\geq0$ in order to write
$A_L[f]=\norm{Gf}^2$ at all. Both sides of
Proposition~\ref{prop:gamma-compression} are positive term by term, at every
$L$, so \eqref{eq:prime-free-target} is not a prerequisite for it.
That inequality remains an open question about $A_L$; no part of the
construction problem now waits on it.

\emph{The price is that the subtracted term is no longer purely arithmetic.}
It carries the contact and the odd half of the pole form as well as every
prime. The trade is a dominating form that can be written down, against a
subtracted form that mixes archimedean and arithmetic data.

The constant is a step function of $L$, jumping by
$\widetilde\kappa_{r_p,d_p}$ at $L=\log p$ from $w_0=-5.37218\ldots$, and its
partial sums first become positive at $p=13$:
\begin{center}
\small
\begin{tabular}{lrrrrrrr}
\toprule
$p$ & $2$ & $3$ & $5$ & $7$ & $11$ & $13$ & $17$\\
$c_L$ after $p$ & $-4.7980$ & $-3.9937$ & $-2.9990$ & $-1.9315$ & $-0.8205$
 & $+0.2933$ & $+1.3994$\\
\bottomrule
\end{tabular}
\end{center}
So $c_L\leq0$ exactly for $L<\log13=2.5649\ldots$, which is the precise form of
the sign change recorded after \eqref{eq:prime-free}; above it the constant
joins the source and only helps.

The presentation does not make the problem smaller, and says so
quantitatively. Since
$\norm\Pi=1-\min\operatorname{spec}(Q_L;K_+)$ and $Q_L$ is nearly degenerate,
the domination is saturated. In a Fourier--Galerkin space of $41$ modes the
largest generalized eigenvalue of $(T_L,K_+)$ is $1-1.9\times10^{-7}$ at $L=1$
and equals $1$ to within $10^{-15}$ at $L=3/2$, $2$ and $3$, as $\norm{\Pi_L}$
does after Corollary~\ref{cor:domination}; the eigenvalues accumulate at one,
with $1,2,4,6,11,24$ of them above $0.999$ at $L=4/5,1,5/4,3/2,2,3$. A
candidate mechanism must therefore produce a \emph{critical} contraction rather
than a strict one, whose spectrum piles up at one on a subspace growing with
$L$; a mechanism yielding a uniform strict contraction would be producing a
false statement. These are one-sided subspace values, as after
Corollary~\ref{cor:domination}: above one they would certify that the
domination fails, below one they certify nothing.
